% Options for packages loaded elsewhere
\PassOptionsToPackage{unicode}{hyperref}
\PassOptionsToPackage{hyphens}{url}
\documentclass[
]{article}
\usepackage{xcolor}
\usepackage[letterpaper,top=2cm,bottom=2cm,left=1.5cm,right=1.5cm]{geometry}
\usepackage{amsmath,amssymb}
\setcounter{secnumdepth}{-\maxdimen} % remove section numbering
\usepackage{iftex}
\ifPDFTeX
  \usepackage[T1]{fontenc}
  \usepackage[utf8]{inputenc}
  \usepackage{textcomp} % provide euro and other symbols
\else % if luatex or xetex
  \usepackage{unicode-math} % this also loads fontspec
  \defaultfontfeatures{Scale=MatchLowercase}
  \defaultfontfeatures[\rmfamily]{Ligatures=TeX,Scale=1}
\fi
\usepackage{lmodern}
\ifPDFTeX\else
  % xetex/luatex font selection
\fi
% Use upquote if available, for straight quotes in verbatim environments
\IfFileExists{upquote.sty}{\usepackage{upquote}}{}
\IfFileExists{microtype.sty}{% use microtype if available
  \usepackage[]{microtype}
  \UseMicrotypeSet[protrusion]{basicmath} % disable protrusion for tt fonts
}{}
\makeatletter
\@ifundefined{KOMAClassName}{% if non-KOMA class
  \IfFileExists{parskip.sty}{%
    \usepackage{parskip}
  }{% else
    \setlength{\parindent}{0pt}
    \setlength{\parskip}{6pt plus 2pt minus 1pt}}
}{% if KOMA class
  \KOMAoptions{parskip=half}}
\makeatother
\setlength{\emergencystretch}{3em} % prevent overfull lines
\providecommand{\tightlist}{%
  \setlength{\itemsep}{0pt}\setlength{\parskip}{0pt}}
\usepackage{bookmark}
\IfFileExists{xurl.sty}{\usepackage{xurl}}{} % add URL line breaks if available
\urlstyle{same}
% fallback for those not using the hyperref driver hyperxmp:
\makeatletter
\@ifundefined{xmpquote}{\newcommand{\xmpquote}[1]{#1}}{}
\makeatother
\hypersetup{
  pdftitle={Cluster Analysis - Lecture Notes},
  pdfauthor={MA2003B},
  hidelinks,
  pdfcreator={LaTeX via pandoc}}

\title{Cluster Analysis - Lecture Notes}
\author{MA2003B}
\date{}

\begin{document}
\maketitle

\textbf{Cluster Analysis}

Lecture Notes

MA2003B - Application of Multivariate Methods in Data Science

\begin{center}\rule{0.5\linewidth}{0.5pt}\end{center}

\textbf{ }

\section{Document Information}

\textbf{ }

\subsection{Credits and Authorship}

\textbf{Content Development:}

\begin{itemize}
\item
  \textbf{AI Assistant:} Claude (Anthropic) - Initial content
  generation, mathematical formulations, and structural organization
\item
  \textbf{Human Reviewer and Course Coordinator:} Juliho Castillo
  Colmenares - Academic review, pedagogical alignment, and quality
  assurance
\end{itemize}

\textbf{ }

\subsection{Roles and Responsibilities}

\textbf{Claude (AI Assistant):}

\begin{itemize}
\item
  Generated comprehensive lecture notes based on cluster analysis
  curriculum
\item
  Developed mathematical formulations and technical explanations
\item
  Created practice questions and detailed answers
\item
  Structured content according to pedagogical best practices
\item
  Aligned material with E06 Cluster Analysis quiz topics
\end{itemize}

\textbf{Human Reviewer (Course Coordinator):}

\begin{itemize}
\item
  Reviewed content for academic accuracy and rigor
\item
  Ensured alignment with MA2003B course learning objectives
\item
  Validated examples and practice questions
\item
  Approved final content for student distribution
\item
  Maintains ongoing responsibility for course materials
\end{itemize}

\textbf{ }

\subsection{Contact Information}

\textbf{Course Coordinator:}

\begin{itemize}
\item
  \textbf{Name:} Juliho Castillo Colmenares
\item
  \textbf{Email:} julihocc@tec
\item
  \textbf{Office:} Tec de Monterrey CCM Office 1540
\item
  \textbf{Office Hours:} Monday-Friday, 9:00 AM - 5:00 PM
\end{itemize}

\textbf{Course Information:}

\begin{itemize}
\item
  \textbf{Course Code:} MA2003B
\item
  \textbf{Course Title:} Application of Multivariate Methods in Data
  Science
\item
  \textbf{Institution:} Tec de Monterrey
\item
  \textbf{Department:} Escuela de Ingeniería y Ciencias
\end{itemize}

\textbf{ }

\subsection{Version Information}

\begin{itemize}
\item
  \textbf{Document Version:} 1.0
\item
  \textbf{Created:} 2025-01-07
\item
  \textbf{Last Updated:} 2025-01-07
\item
  \textbf{Status:} Active - Current Semester
\end{itemize}

\textbf{ }

\subsection{Usage and Distribution}

These lecture notes are intended for students enrolled in MA2003B. The
material is designed to prepare students for the Cluster Analysis
evaluation (E06) and provide comprehensive understanding of unsupervised
learning methods.

\textbf{Academic Integrity:} Students are expected to use these notes as
a study resource in accordance with institutional academic integrity
policies.

\begin{center}\rule{0.5\linewidth}{0.5pt}\end{center}

\begin{center}\rule{0.5\linewidth}{0.5pt}\end{center}

\textbf{ }

\section{Introduction to Cluster Analysis}

\textbf{ }

\subsection{What is Cluster Analysis?}

\textbf{Cluster Analysis} is an exploratory data analysis technique used
to discover natural groupings (clusters) in data \textbf{without
predefined categories}. Unlike discriminant analysis, which classifies
observations into known groups, cluster analysis seeks to identify
previously unknown structure in the data.

\textbf{Info: } \textbf{Key Distinction:} Cluster analysis is
\textbf{unsupervised learning} - we don't know the ``true'' groups
beforehand. We let the data reveal its own structure.

\textbf{ }

\subsection{Primary Objectives}

\begin{enumerate}
\item
  \textbf{Grouping:} Partition observations into homogeneous groups
\item
  \textbf{Data reduction:} Simplify large datasets by representing
  groups
\item
  \textbf{Pattern discovery:} Identify natural structure in data
\item
  \textbf{Hypothesis generation:} Suggest new classifications for
  further study
\end{enumerate}

\textbf{ }

\subsection{Common Applications}

\begin{itemize}
\item
  \textbf{Marketing:} Customer segmentation for targeted campaigns
\item
  \textbf{Biology:} Taxonomy and species classification
\item
  \textbf{Medicine:} Disease subtype identification
\item
  \textbf{Social Sciences:} Community detection in networks
\item
  \textbf{Image Processing:} Image segmentation and compression
\end{itemize}

\begin{center}\rule{0.5\linewidth}{0.5pt}\end{center}

\textbf{ }

\section{Distance and Similarity Measures}

\textbf{ }

\subsection{Why Distance Matters}

Clustering relies on measuring how ``close'' or ``similar'' observations
are to each other. The choice of distance metric fundamentally affects
the clustering results.

\textbf{ }

\subsection{Common Distance Measures}

\subsubsection{Euclidean Distance (L2 Norm)}

The most commonly used distance measure, representing straight-line
distance in n-dimensional space.

\[d(x,y) = \sqrt{\sum_{i = 1}^{p}\left( x_{i} - y_{i} \right)^{2}}\]

\textbf{Properties:}

\begin{itemize}
\item
  Sensitive to scale differences between variables
\item
  Assumes equal importance of all dimensions
\item
  Works best when variables are continuous and measured on similar
  scales
\end{itemize}

\textbf{Example: } For points \(x = (2,3)\) and \(y = (5,7)\):

\(d(x,y) = \sqrt{(5 - 2)^{2} + (7 - 3)^{2}} = \sqrt{9 + 16} = 5\)

\subsubsection{Manhattan Distance (L1 Norm)}

Also called ``city block'' distance, it measures distance as the sum of
absolute differences.

\[d(x,y) = \sum_{i = 1}^{p}|x_{i} - y_{i}|\]

\textbf{When to use:}

\begin{itemize}
\item
  Data contains outliers or extreme values (more robust than Euclidean)
\item
  Variables represent counts or discrete quantities
\item
  Working in high-dimensional spaces
\end{itemize}

\textbf{Example: } For the same points \(x = (2,3)\) and \(y = (5,7)\):

\(d(x,y) = |5 - 2| + |7 - 3| = 3 + 4 = 7\)

\subsubsection{Other Distance Measures}

\textbf{Cosine Similarity:} Measures angle between vectors (useful for
text data)

\[\text{ similarity}(x,y) = \frac{x \cdot y}{\| x\| \cdot \| y\|}\]

\textbf{Correlation-based Distance:} \(d(x,y) = 1 - r(x,y)\)

Used when pattern similarity is more important than magnitude.

\textbf{Warning: } \textbf{Critical:} Always standardize variables
measured on different scales before computing distances! Variables with
larger scales will dominate the distance calculation otherwise.

\textbf{ }

\subsection{Standardization}

Before clustering, transform variables to have mean 0 and standard
deviation 1:

\[z_{i} = \frac{x_{i} - \mu}{\sigma}\]

This ensures all variables contribute equally to distance calculations.

\begin{center}\rule{0.5\linewidth}{0.5pt}\end{center}

\textbf{ }

\section{Hierarchical Clustering}

\textbf{ }

\subsection{Overview}

Hierarchical clustering builds a tree-like structure (dendrogram)
showing nested clusters at different levels of granularity. It doesn't
require specifying the number of clusters in advance.

\textbf{ }

\subsection{Types of Hierarchical Clustering}

\subsubsection{Agglomerative (Bottom-Up)}

\begin{itemize}
\item
  \textbf{Start:} Each observation is its own cluster
\item
  \textbf{Process:} Iteratively merge the two closest clusters
\item
  \textbf{End:} All observations in one cluster
\item
  \textbf{Most common approach}
\end{itemize}

\subsubsection{Divisive (Top-Down)}

\begin{itemize}
\item
  \textbf{Start:} All observations in one cluster
\item
  \textbf{Process:} Iteratively split the most heterogeneous cluster
\item
  \textbf{End:} Each observation is its own cluster
\item
  \textbf{Less common, computationally intensive}
\end{itemize}

\textbf{ }

\subsection{Linkage Methods}

Linkage methods define how to measure distance between \textbf{clusters}
(not just individual points).

\subsubsection{Single Linkage (Nearest Neighbor)}

Distance = \textbf{minimum} distance between any two points in the
clusters

\[d\left( C_{1},C_{2} \right) = \min\limits_{x \in C_{1},y \in C_{2}}d(x,y)\]

\textbf{Advantages:}

\begin{itemize}
\item
  Can detect elongated, non-spherical clusters
\item
  Handles irregular cluster shapes
\end{itemize}

\textbf{Disadvantages:}

\begin{itemize}
\item
  Susceptible to ``chaining effect'' - clusters connect via single
  intermediate points
\item
  Sensitive to noise and outliers
\end{itemize}

\subsubsection{Complete Linkage (Farthest Neighbor)}

Distance = \textbf{maximum} distance between any two points in the
clusters

\[d\left( C_{1},C_{2} \right) = \max\limits_{x \in C_{1},y \in C_{2}}d(x,y)\]

\textbf{Advantages:}

\begin{itemize}
\item
  Avoids single linkage's chaining effect (does not merge clusters based
  on one close pair while ignoring the rest)
\item
  Tends to produce compact, spherical clusters
\end{itemize}

\textbf{Disadvantages:}

\begin{itemize}
\item
  May break large clusters
\item
  Biased toward finding clusters of similar size
\item
  Its \textbf{maximum}-pairwise-distance criterion is itself driven by
  whichever two points are farthest apart, so a single outlier in either
  cluster can dominate the linkage distance -- complete linkage is not
  generally ``robust to outliers,'' only less prone than single linkage
  to the specific failure mode of chaining
\end{itemize}

\subsubsection{Average Linkage}

Distance = \textbf{average} of all pairwise distances between clusters

\[d\left( C_{1},C_{2} \right) = \frac{1}{n_{1}n_{2}}\sum_{x \in C_{1}}\sum_{y \in C_{2}}d(x,y)\]

\textbf{Advantages:}

\begin{itemize}
\item
  Compromise between single and complete linkage
\item
  Less affected by outliers than single linkage
\item
  Generally produces good results
\end{itemize}

\subsubsection{Ward's Method}

Minimizes the \textbf{total within-cluster sum of squares} (variance),
\textbf{provided the dissimilarity is squared Euclidean distance} --
this variance-minimizing interpretation does not hold for other distance
metrics (e.g. Manhattan, cosine), even though Ward's linkage update rule
can technically be run on any dissimilarity matrix.

At each step, merge the two clusters that result in the smallest
increase in total within-cluster variance.

\[\text{ ESS } = \sum_{i = 1}^{k}\sum_{x \in C_{i}}\| x - \mu_{i}\|_{2}^{2}\]

\textbf{Advantages:}

\begin{itemize}
\item
  Tends to create compact, equal-sized clusters
\item
  Minimizes variance within clusters
\item
  Often produces the best results in practice
\end{itemize}

\textbf{Disadvantages:}

\begin{itemize}
\item
  Assumes spherical clusters
\item
  Sensitive to outliers
\end{itemize}

\textbf{Tip: } \textbf{Practical Recommendation:} Ward's method is often
the best choice for exploratory analysis. Use average linkage as an
alternative. Avoid single linkage unless you specifically need to detect
elongated clusters.

\textbf{ }

\subsection{Dendrograms}

A dendrogram visualizes the hierarchical clustering process as a tree
diagram.

\textbf{Reading a Dendrogram:}

\begin{itemize}
\item
  Horizontal axis: Individual observations or clusters
\item
  Vertical axis: Distance or dissimilarity at which clusters merge
\item
  Height of branches: Indicates the distance between merged clusters
\end{itemize}

\textbf{Determining Number of Clusters:} Look for large vertical gaps
(jumps in distance) - cut the dendrogram where there's a substantial
increase in fusion distance.

\textbf{Info: } \textbf{Cutting the Dendrogram:} Draw a horizontal line
across the dendrogram. The number of vertical lines it crosses equals
the number of clusters at that level.

\textbf{ }

\subsection{Chaining Effect}

The ``chaining effect'' occurs when clusters form long, elongated chains
rather than compact groups. This happens with single linkage when
observations connect via intermediate points.

\textbf{Example scenario:} Observations A-B-C-D form a chain where each
is close to its neighbor, but A and D are far apart. Single linkage
would group them together.

\begin{center}\rule{0.5\linewidth}{0.5pt}\end{center}

\textbf{ }

\section{Non-Hierarchical Clustering}

\textbf{ }

\subsection{K-Means Clustering}

The most popular non-hierarchical method. Partitions data into k
clusters by minimizing within-cluster variance.

\subsubsection{Algorithm}

\begin{enumerate}
\item
  \textbf{Initialize:} Randomly select k observations as initial cluster
  centers (centroids)
\item
  \textbf{Assignment:} Assign each observation to the nearest centroid
\item
  \textbf{Update:} Recalculate centroids as the mean of all points in
  each cluster
\item
  \textbf{Repeat:} Steps 2-3 until assignments no longer change
  (convergence)
\end{enumerate}

\subsubsection{Mathematical Objective}

Minimize the total within-cluster sum of squares (WCSS):

\[\min\sum_{i = 1}^{k}\sum_{x \in C_{i}}\| x - \mu_{i}\|^{2}\]

where \(\mu_{i}\) is the centroid of cluster \(C_{i}\).

\subsubsection{Convergence}

The algorithm \textbf{always converges} because:

\begin{itemize}
\item
  Each iteration reduces or maintains WCSS
\item
  There are finitely many possible partitions
\item
  Typically converges in 10-30 iterations
\end{itemize}

Convergence is reached when cluster assignments no longer change between
iterations.

\subsubsection{Advantages}

\begin{itemize}
\item
  \textbf{Fast:} Computational complexity O(n × k × p × iterations)
\item
  \textbf{Scalable:} Works well with large datasets
\item
  \textbf{Simple:} Easy to understand and implement
\item
  \textbf{Efficient:} Generally converges quickly
\end{itemize}

\subsubsection{Limitations}

\begin{enumerate}
\item
  \textbf{Requires specifying k in advance} - must know or guess number
  of clusters
\item
  \textbf{Sensitive to initialization} - different starting points may
  yield different results
\item
  \textbf{Assumes spherical clusters} - struggles with elongated or
  irregular shapes
\item
  \textbf{Sensitive to outliers} - outliers can distort centroids
\item
  \textbf{Assumes equal cluster sizes} - tends to create similar-sized
  clusters
\end{enumerate}

\subsubsection{K-Means++ Initialization}

An improved initialization method that spreads out initial centroids:

\begin{enumerate}
\item
  Choose first centroid randomly
\item
  For each subsequent centroid, choose a point with probability
  proportional to its squared distance from nearest existing centroid
\item
  Repeat until k centroids selected
\end{enumerate}

\textbf{Benefit:} Significantly reduces the risk of poor initialization
and typically produces better results.

\textbf{ }

\subsection{K-Medoids (PAM - Partitioning Around Medoids)}

Similar to k-means but uses actual data points as cluster centers
(medoids) instead of computed means.

\subsubsection{Key Difference from K-Means}

\begin{itemize}
\item
  \textbf{K-means:} Centers are computed means (may not be actual data
  points)
\item
  \textbf{K-medoids:} Centers are actual data points from the dataset
\end{itemize}

\subsubsection{Advantages}

\begin{itemize}
\item
  \textbf{More robust to outliers than k-means} - an extreme point can
  still become (or shift) a medoid, but a medoid is constrained to be an
  actual observation, so it cannot be dragged arbitrarily far the way a
  computed mean can
\item
  \textbf{Works with any distance metric} - not limited to Euclidean
  distance
\item
  \textbf{Interpretable centers} - medoids are actual observations
\end{itemize}

\subsubsection{Disadvantages}

\begin{itemize}
\item
  \textbf{Slower:} Higher computational complexity than k-means
\item
  \textbf{Less common:} Not as widely implemented in software
\end{itemize}

\textbf{Tip: } \textbf{When to use K-medoids:} Choose k-medoids when
your data contains outliers or when you need cluster centers to be
actual observations (e.g., choosing representative customers).

\begin{center}\rule{0.5\linewidth}{0.5pt}\end{center}

\textbf{ }

\section{Determining Optimal Number of Clusters}

\textbf{ }

\subsection{The Fundamental Challenge}

Unlike supervised learning, we don't have a ``ground truth'' for the
correct number of clusters. We must use heuristics and domain knowledge.

\textbf{ }

\subsection{Elbow Method}

Plot the total within-cluster sum of squares (WCSS) against the number
of clusters k.

\textbf{Procedure:}

\begin{enumerate}
\item
  Run clustering for k = 1, 2, 3, \ldots, K\_max
\item
  Calculate WCSS for each k
\item
  Plot WCSS vs. k
\item
  Look for an ``elbow'' - the point where adding more clusters yields
  diminishing returns
\end{enumerate}

\[\text{ WCSS}(k) = \sum_{i = 1}^{k}\sum_{x \in C_{i}}\| x - \mu_{i}\|^{2}\]

\textbf{Interpretation:}

\begin{itemize}
\item
  WCSS always decreases as k increases
\item
  The ``elbow'' indicates where additional clusters don't substantially
  improve fit
\item
  Choose k at the elbow point
\end{itemize}

\textbf{Warning: } \textbf{Limitation:} The elbow is not always clear -
sometimes the curve is smooth with no obvious bend.

\textbf{ }

\subsection{Silhouette Analysis}

Measures how well each point fits within its assigned cluster compared
to other clusters.

\subsubsection{Silhouette Coefficient}

For each observation i:

\[s(i) = \frac{b(i) - a(i)}{\max(a(i),b(i))}\]

where:

\begin{itemize}
\item
  \(a(i)\) = average distance to other points in the same cluster
  (cohesion)
\item
  \(b(i)\) = average distance to points in the nearest neighboring
  cluster (separation)
\end{itemize}

\textbf{Range:} \(s(i) \in \lbrack - 1, + 1\rbrack\)

\textbf{Interpretation:}

\begin{itemize}
\item
  \(s(i) \approx + 1\): Point is well-matched to its cluster (far from
  neighbors)
\item
  \(s(i) \approx 0\): Point is on the border between clusters
\item
  \(s(i) \approx - 1\): Point may be assigned to the wrong cluster
\end{itemize}

\subsubsection{Average Silhouette Width}

Calculate the mean silhouette coefficient across all observations:

\[\overline{s} = \frac{1}{n}\sum_{i = 1}^{n}s(i)\]

\textbf{Usage:} Run clustering for different values of k and choose the
k that maximizes the average silhouette width.

\textbf{Info: } \textbf{Silhouette Plot:} Create a bar chart showing
silhouette coefficients for each observation, grouped by cluster. Wide,
tall bars indicate good clustering.

\textbf{ }

\subsection{Other Methods}

\subsubsection{Gap Statistic}

Compares WCSS to its expected value under a null reference distribution
(random data):

\[\operatorname{Gap}(k) = {\mathbb{E}}^{\ast}\left\lbrack \log(\text{WCSS}_{k}) \right\rbrack - \log(\text{WCSS}_{k})\]

where the expectation is estimated by averaging over many reference
datasets sampled uniformly over the observed data's range.

\textbf{Selection rule:} Simply picking the \(k\) with the largest
\(\operatorname{Gap}(k)\) ignores the Monte Carlo simulation error in
estimating the expectation. The standard rule (Tibshirani, Walther \&
Hastie, 2001) instead uses the \textbf{one-standard-error} criterion:
letting \(s_{k}\) be the estimated standard error of
\(\operatorname{Gap}(k)\), choose the \textbf{smallest} \(k\) such that
\[\operatorname{Gap}(k) \geq \operatorname{Gap}(k + 1) - s_{k + 1}\]
i.e. the smallest \(k\) whose gap is not significantly exceeded by the
next larger \(k\).

\subsubsection{Davies-Bouldin Index}

Measures the ratio of within-cluster dispersion to between-cluster
separation.

\[\text{ DB } = \frac{1}{k}\sum_{i = 1}^{k}\max\limits_{j \neq i}\left( \frac{\sigma_{i} + \sigma_{j}}{d\left( c_{i},c_{j} \right)} \right)\]

\textbf{Lower values indicate better clustering} - compact clusters that
are far apart.

\subsubsection{Domain Knowledge}

Often the most important factor:

\begin{itemize}
\item
  Business requirements (e.g., ``we need 3-5 customer segments for our
  marketing capacity'')
\item
  Interpretability (fewer clusters are easier to explain)
\item
  Actionability (can we act differently for each cluster?)
\end{itemize}

\begin{center}\rule{0.5\linewidth}{0.5pt}\end{center}

\textbf{ }

\section{Validation and Quality Assessment}

\textbf{ }

\subsection{Internal Validation}

Measures clustering quality using only the data itself (no external
labels).

\subsubsection{Within-Cluster Sum of Squares (WCSS)}

\[\text{ WCSS } = \sum_{i = 1}^{k}\sum_{x \in C_{i}}\| x - \mu_{i}\|^{2}\]

\textbf{Lower is better} - indicates tight, compact clusters.

\textbf{Limitation:} Always decreases with more clusters, reaching 0
when k = n.

\subsubsection{Between-Cluster Sum of Squares (BCSS)}

\[\text{ BCSS } = \sum_{i = 1}^{k}n_{i}\|\mu_{i} - \mu\|^{2}\]

where \(\mu\) is the overall data mean.

\textbf{Higher is better} - indicates well-separated clusters.

\subsubsection{Silhouette Coefficient}

Already discussed - ranges from -1 to +1, higher is better.

\subsubsection{Davies-Bouldin Index}

\[\text{ DB } = \frac{1}{k}\sum_{i = 1}^{k}\max\limits_{j \neq i}\left( \frac{\sigma_{i} + \sigma_{j}}{d\left( c_{i},c_{j} \right)} \right)\]

\textbf{Lower is better} - indicates compact, well-separated clusters.

\subsubsection{Dunn Index}

\[\text{ Dunn } = \frac{\min\limits_{i \neq j}d\left( C_{i},C_{j} \right)}{\max\limits_{k}\text{ diameter}\left( C_{k} \right)}\]

Ratio of minimum inter-cluster distance to maximum cluster diameter.

\textbf{Higher is better} - indicates well-separated, compact clusters.

\textbf{ }

\subsection{External Validation}

When ``true'' labels are available (e.g., for method comparison).

\subsubsection{Adjusted Rand Index (ARI)}

Measures agreement between clustering results and true labels, adjusted
for chance.

\textbf{Range:} {[}-1, +1{]}

\begin{itemize}
\item
  1 = perfect agreement
\item
  0 = random clustering
\item
  Negative = worse than random
\end{itemize}

\subsubsection{Normalized Mutual Information (NMI)}

Measures information shared between clustering and true labels.

\textbf{Range:} {[}0, 1{]}, higher is better.

\begin{center}\rule{0.5\linewidth}{0.5pt}\end{center}

\textbf{ }

\section{Practical Considerations}

\textbf{ }

\subsection{Curse of Dimensionality}

As the number of dimensions (variables) increases:

\begin{enumerate}
\item
  \textbf{Distance becomes less meaningful} - all points appear
  equidistant in high dimensions
\item
  \textbf{Data becomes sparse} - observations spread out, making
  clusters harder to find
\item
  \textbf{Computational cost increases} - both time and memory
  requirements grow
\end{enumerate}

\textbf{Solutions:}

\begin{itemize}
\item
  \textbf{Dimensionality reduction:} Use PCA or feature selection before
  clustering
\item
  \textbf{Feature selection:} Choose only relevant variables
\item
  \textbf{Specialized methods:} Use algorithms designed for high
  dimensions
\end{itemize}

\textbf{Warning: } \textbf{Rule of Thumb:} If p (dimensions) is large
relative to n (observations), consider dimensionality reduction before
clustering.

\textbf{ }

\subsection{Standardization and Scaling}

\subsubsection{Why Standardize?}

Variables measured on different scales will dominate distance
calculations.

\textbf{Example:} Age (20-80) vs. Income (20,000-200,000)

\begin{itemize}
\item
  Without standardization, income dominates
\item
  Euclidean distance primarily reflects income differences
\end{itemize}

\subsubsection{When to Standardize}

\textbf{Always standardize when:}

\begin{itemize}
\item
  Variables have different units (e.g., kg, cm, years)
\item
  Variables have different scales (e.g., 0-1 vs. 0-1000)
\item
  You want all variables to contribute equally
\end{itemize}

\textbf{Consider not standardizing when:}

\begin{itemize}
\item
  All variables have the same units and scale
\item
  Scale differences are meaningful (e.g., different importance)
\item
  Using correlation-based distances
\end{itemize}

\subsubsection{Standardization Methods}

\textbf{Z-score standardization:} \[z = \frac{x - \mu}{\sigma}\]

\textbf{Min-max scaling:} \[x' = \frac{x - \min(x)}{\max(x) - \min(x)}\]

\begin{center}\rule{0.5\linewidth}{0.5pt}\end{center}

\textbf{ }

\subsection{Handling Mixed Data Types}

\subsubsection{Continuous + Categorical Variables}

\textbf{Option 1: Gower Distance}

Combines different distance measures for different variable types:

\[d_{\text{Gower}(i,j)} = \frac{\sum_{k = 1}^{p}\delta_{ijk}d_{ijk}}{\sum_{k = 1}^{p}\delta_{ijk}}\]

where:

\begin{itemize}
\item
  For continuous:
  \(d_{ijk} = |x_{ik} - x_{jk}\frac{|}{\text{ range}_{k}}\)
\item
  For categorical: \(d_{ijk} = 0\) if same, 1 if different
\item
  \(\delta_{ijk} = 1\) if comparison is valid, 0 if missing
\end{itemize}

\textbf{Option 2: Dummy Encoding}

Convert categorical variables to binary dummy variables, then use
standard distance measures.

\textbf{Caution:} This changes the geometry of the data.

\subsubsection{Pure Categorical Data}

Use specialized distance measures:

\begin{itemize}
\item
  \textbf{Matching coefficient:} Proportion of matching categories
\item
  \textbf{Jaccard distance:} For binary data
\item
  \textbf{Hamming distance:} Number of position differences
\end{itemize}

\textbf{ }

\subsection{Outliers and Robustness}

\subsubsection{Impact of Outliers}

\begin{itemize}
\item
  \textbf{K-means:} Highly sensitive - outliers pull centroids away from
  dense regions
\item
  \textbf{Hierarchical (Ward's):} Sensitive - outliers inflate variance
\item
  \textbf{Single linkage:} Moderately sensitive - may create singleton
  clusters
\item
  \textbf{K-medoids:} More robust than k-means - medoids are constrained
  to be actual observations, but an outlier can still be selected as (or
  influence) a medoid
\end{itemize}

\subsubsection{Strategies for Handling Outliers}

\begin{enumerate}
\item
  \textbf{Pre-processing:} Detect and remove outliers before clustering
\item
  \textbf{Robust methods:} Use k-medoids or other robust algorithms
\item
  \textbf{Outlier clusters:} Accept that some clusters may be outliers
\item
  \textbf{Trimming:} Use trimmed versions of distance measures
\end{enumerate}

\textbf{ }

\subsection{Computational Considerations}

\subsubsection{Complexity}

\begin{itemize}
\item
  \textbf{K-means:} O(n × k × p × iterations) - very efficient
\item
  \textbf{Hierarchical:} O(\(n^{2}\) × p) for distance matrix,
  O(\(n^{3}\)) for some linkages
\item
  \textbf{K-medoids:} O(\(k(n - k)^{2}\) × p) - slower than k-means
\end{itemize}

\subsubsection{Scalability Recommendations}

\begin{itemize}
\item
  \textbf{Small datasets (n \textless{} 1,000):} Any method works
\item
  \textbf{Medium datasets (n = 1,000-10,000):} K-means or k-medoids
\item
  \textbf{Large datasets (n \textgreater{} 10,000):} K-means with
  sampling, or specialized algorithms
\item
  \textbf{Very large datasets (n \textgreater{} 1M):} Mini-batch
  k-means, approximate methods
\end{itemize}

\begin{center}\rule{0.5\linewidth}{0.5pt}\end{center}

\textbf{ }

\section{Cluster Analysis Workflow}

\textbf{ }

\subsection{Step-by-Step Procedure}

\subsubsection{1. Problem Definition}

\begin{itemize}
\item
  Define the objective: What questions are we trying to answer?
\item
  Determine what makes a ``good'' cluster in your context
\item
  Consider how results will be used
\end{itemize}

\subsubsection{2. Variable Selection}

\begin{itemize}
\item
  Choose relevant variables (domain knowledge)
\item
  Remove redundant variables (high correlation)
\item
  Consider creating derived features
\item
  Decide whether to include or exclude certain variables
\end{itemize}

\subsubsection{3. Data Preprocessing}

\textbf{Handle Missing Values:}

\begin{itemize}
\item
  Imputation
\item
  Deletion
\item
  Special clustering methods for incomplete data
\end{itemize}

\textbf{Deal with Outliers:}

\begin{itemize}
\item
  Detect using box plots, z-scores, Mahalanobis distance
\item
  Decide: remove, transform, or use robust methods
\end{itemize}

\textbf{Transform Variables:}

\begin{itemize}
\item
  Log transform for skewed distributions
\item
  Square root for count data
\item
  Box-Cox transformations
\end{itemize}

\textbf{Standardize:}

\begin{itemize}
\item
  Z-score standardization if variables on different scales
\item
  Consider domain knowledge about variable importance
\end{itemize}

\subsubsection{4. Choose Clustering Method}

\textbf{Hierarchical if:}

\begin{itemize}
\item
  Small to medium dataset
\item
  Want to explore different numbers of clusters
\item
  Need hierarchical structure (taxonomy)
\item
  Don't know k in advance
\end{itemize}

\textbf{K-means if:}

\begin{itemize}
\item
  Large dataset
\item
  Approximately know the number of clusters
\item
  Need speed and efficiency
\item
  Clusters are roughly spherical
\end{itemize}

\textbf{K-medoids if:}

\begin{itemize}
\item
  Outliers present
\item
  Need cluster centers to be actual observations
\item
  Can afford extra computation
\end{itemize}

\textbf{Other methods if:}

\begin{itemize}
\item
  Non-spherical clusters (DBSCAN)
\item
  Mixed data types (k-prototypes)
\item
  High dimensions (subspace clustering)
\end{itemize}

\subsubsection{5. Determine Number of Clusters}

\begin{itemize}
\item
  Use elbow method
\item
  Calculate silhouette coefficients
\item
  Try multiple values of k
\item
  Consider business constraints
\item
  Validate interpretability
\end{itemize}

\subsubsection{6. Run Clustering}

\begin{itemize}
\item
  For k-means: run multiple times with different initializations
\item
  For hierarchical: create dendrogram and examine
\item
  Save cluster assignments and centers
\end{itemize}

\subsubsection{7. Validate Results}

\textbf{Internal validation:}

\begin{itemize}
\item
  Check silhouette coefficients
\item
  Examine within/between cluster variance
\item
  Look for negative silhouette values
\end{itemize}

\textbf{Stability:}

\begin{itemize}
\item
  Re-run with bootstrap samples
\item
  Check if clusters are consistent
\end{itemize}

\textbf{Interpretability:}

\begin{itemize}
\item
  Can you describe each cluster clearly?
\item
  Do clusters make business sense?
\end{itemize}

\subsubsection{8. Interpret and Profile Clusters}

\begin{itemize}
\item
  Examine cluster centers/medoids
\item
  Calculate cluster means for each variable
\item
  Create visualizations (parallel coordinates, scatter plots)
\item
  Name/label clusters based on characteristics
\item
  Identify distinguishing features of each cluster
\end{itemize}

\subsubsection{9. Validation and Iteration}

\begin{itemize}
\item
  Test on new data if possible
\item
  Refine based on domain expert feedback
\item
  Consider alternative numbers of clusters
\item
  Try different methods or distance measures
\end{itemize}

\subsubsection{10. Report and Act}

\begin{itemize}
\item
  Document methodology and rationale
\item
  Present cluster profiles clearly
\item
  Provide actionable insights
\item
  Discuss limitations and next steps
\end{itemize}

\begin{center}\rule{0.5\linewidth}{0.5pt}\end{center}

\textbf{ }

\section{Advanced Topics}

\textbf{ }

\subsection{Density-Based Clustering}

\subsubsection{DBSCAN (Density-Based Spatial Clustering)}

Unlike k-means, DBSCAN can find clusters of arbitrary shape and identify
outliers.

\textbf{Key Parameters:}

\begin{itemize}
\item
  \(\varepsilon\): neighborhood radius
\item
  MinPts: minimum points to form dense region
\end{itemize}

\textbf{Advantages:}

\begin{itemize}
\item
  Finds non-spherical clusters
\item
  Robust to outliers (marks them as noise)
\item
  No need to specify number of clusters
\end{itemize}

\textbf{Disadvantages:}

\begin{itemize}
\item
  Struggles with varying densities
\item
  Sensitive to parameters
\item
  Higher computational cost
\end{itemize}

\textbf{ }

\subsection{Fuzzy Clustering}

Each observation has a degree of membership in each cluster (soft
assignment) rather than belonging to exactly one cluster (hard
assignment).

\textbf{Fuzzy C-means:}

\begin{itemize}
\item
  Membership values range from 0 to 1
\item
  Memberships for each point sum to 1
\item
  Useful when boundaries between clusters are unclear
\end{itemize}

\textbf{ }

\subsection{Model-Based Clustering}

Assumes data comes from a mixture of probability distributions
(typically Gaussian).

\textbf{Gaussian Mixture Models (GMM):}

\begin{itemize}
\item
  Each cluster is a Gaussian distribution
\item
  Use EM algorithm to estimate parameters
\item
  Provides probabilistic cluster assignments
\item
  Can estimate optimal k using BIC or AIC
\end{itemize}

\textbf{ }

\subsection{Subspace and Projected Clustering}

For high-dimensional data, different clusters may exist in different
subspaces (subsets of variables).

\textbf{Approaches:}

\begin{itemize}
\item
  Find clusters in different variable subsets
\item
  Identify relevant dimensions for each cluster
\item
  Useful when curse of dimensionality is severe
\end{itemize}

\begin{center}\rule{0.5\linewidth}{0.5pt}\end{center}

\textbf{ }

\section{Common Pitfalls and Best Practices}

\textbf{ }

\subsection{Common Mistakes to Avoid}

\textbf{Warning: } \textbf{Pitfall 1: Not standardizing variables}

Variables with larger scales will dominate distance calculations,
leading to meaningless clusters based primarily on scale rather than
pattern.

\textbf{Warning: } \textbf{Pitfall 2: Using k-means with non-spherical
clusters}

K-means assumes spherical clusters and will force non-spherical data
into inappropriate spherical groups.

\textbf{Warning: } \textbf{Pitfall 3: Ignoring outliers}

Outliers can severely distort clustering results, especially with
k-means and Ward's method.

\textbf{Warning: } \textbf{Pitfall 4: Over-interpreting results}

Clustering will always find some structure, even in random data.
Validate that clusters are meaningful and stable.

\textbf{Warning: } \textbf{Pitfall 5: Using too many variables}

High dimensionality makes clustering difficult. Focus on relevant,
non-redundant variables.

\textbf{ }

\subsection{Best Practices}

\textbf{Tip: } \textbf{Practice 1: Try multiple methods}

Compare hierarchical, k-means, and other approaches. Robust clusters
will appear across methods.

\textbf{Tip: } \textbf{Practice 2: Validate stability}

Run clustering multiple times with different initializations or
bootstrap samples. Good clusters should be consistent.

\textbf{Tip: } \textbf{Practice 3: Visualize extensively}

Use scatter plots, parallel coordinates, heatmaps, and dendrograms to
understand cluster structure.

\textbf{Tip: } \textbf{Practice 4: Use domain knowledge}

Statistical metrics are important, but clusters must make practical
sense in your domain.

\textbf{Tip: } \textbf{Practice 5: Document decisions}

Record why you chose certain methods, parameters, and number of
clusters. Clustering involves many subjective choices.

\begin{center}\rule{0.5\linewidth}{0.5pt}\end{center}

\textbf{ }

\section{Summary and Key Takeaways}

\textbf{ }

\subsection{Fundamental Concepts}

\begin{enumerate}
\item
  \textbf{Cluster analysis discovers natural groupings} in data without
  predefined categories
\item
  \textbf{Distance measures are crucial} - choice affects results
  fundamentally:

  \begin{itemize}
  \item
    Euclidean for continuous, similar scales
  \item
    Manhattan for robustness to outliers
  \item
    Correlation for pattern similarity
  \end{itemize}
\item
  \textbf{Standardization is essential} when variables have different
  scales
\end{enumerate}

\textbf{ }

\subsection{Clustering Methods}

\textbf{Hierarchical Clustering:}

\begin{itemize}
\item
  Creates tree structure (dendrogram)
\item
  Don't need to specify k in advance
\item
  Linkage method matters:

  \begin{itemize}
  \item
    Single: can find elongated clusters, prone to chaining
  \item
    Complete: compact clusters, robust to outliers
  \item
    Average: good compromise
  \item
    Ward's: minimizes variance, often best results
  \end{itemize}
\end{itemize}

\textbf{K-means:}

\begin{itemize}
\item
  Fast, scalable, widely used
\item
  Requires specifying k
\item
  Assumes spherical clusters
\item
  Sensitive to initialization and outliers
\item
  Use k-means++ for better initialization
\end{itemize}

\textbf{K-medoids:}

\begin{itemize}
\item
  More robust to outliers than k-means
\item
  Centers are actual data points
\item
  Slower than k-means
\end{itemize}

\textbf{ }

\subsection{Determining Optimal k}

\begin{enumerate}
\item
  \textbf{Elbow method:} Look for bend in WCSS vs. k plot
\item
  \textbf{Silhouette analysis:} Choose k that maximizes average
  silhouette width
\item
  \textbf{Domain knowledge:} Consider business requirements and
  interpretability
\item
  \textbf{Try multiple values:} Compare and validate
\end{enumerate}

\textbf{ }

\subsection{Validation}

\begin{itemize}
\item
  \textbf{Silhouette coefficient:} +1 = perfect, 0 = borderline, -1 =
  wrong cluster
\item
  \textbf{Davies-Bouldin Index:} Lower is better
\item
  \textbf{Stability:} Rerun with different seeds/samples
\item
  \textbf{Interpretability:} Can you explain and use the clusters?
\end{itemize}

\textbf{ }

\subsection{Practical Workflow}

\begin{enumerate}
\item
  Define objective and select variables
\item
  Preprocess: handle missing values, outliers, transform, standardize
\item
  Choose method based on data characteristics and goals
\item
  Determine k using multiple criteria
\item
  Run clustering (multiple times if k-means)
\item
  Validate and interpret results
\item
  Refine and iterate as needed
\end{enumerate}

\textbf{ }

\subsection{Critical Considerations}

\begin{itemize}
\item
  \textbf{Curse of dimensionality:} Reduce dimensions if p is large
\item
  \textbf{Outliers:} Detect and handle appropriately
\item
  \textbf{Mixed data types:} Use specialized distances or preprocessing
\item
  \textbf{Scale matters:} Always consider whether to standardize
\item
  \textbf{No ``true'' answer:} Clustering is exploratory; validate and
  iterate
\end{itemize}

\begin{center}\rule{0.5\linewidth}{0.5pt}\end{center}

\textbf{ }

\section{Practice Questions}

Test your understanding with these questions based on the lecture
material:

\textbf{ }

\subsection{Conceptual Questions}

\begin{enumerate}
\item
  What is the fundamental difference between cluster analysis and
  discriminant analysis?
\item
  Why is standardization important before clustering? Give an example
  where lack of standardization would cause problems.
\item
  Explain the ``chaining effect'' in hierarchical clustering. Which
  linkage method is most susceptible to it?
\item
  What are the main advantages and disadvantages of k-means clustering?
\item
  How does the silhouette coefficient help determine if a point is
  well-clustered?
\end{enumerate}

\textbf{ }

\subsection{Application Questions}

\begin{enumerate}
\setcounter{enumi}{5}
\item
  You have customer data with variables: age (years), income (dollars),
  purchases (count), and loyalty\_score (0-100). Should you standardize
  before clustering? Why?
\item
  Your elbow plot shows a smooth curve with no clear ``elbow''. What
  should you do?
\item
  You run k-means with k=3 five times and get different results each
  time. What might explain this, and how can you address it?
\item
  Your dataset has 50,000 observations and 20 variables. Which
  clustering method would you recommend and why?
\item
  After clustering customers, you find one cluster has 5,000 members and
  another has only 50. Is this a problem? What might it indicate?
\end{enumerate}

\textbf{ }

\subsection{Interpretation Questions}

\begin{enumerate}
\setcounter{enumi}{10}
\item
  What does it mean if a point has a silhouette coefficient of -0.3?
\item
  You get a Davies-Bouldin Index of 0.5 for k=3 and 0.8 for k=5. Which
  is better?
\item
  In a dendrogram, you see a large vertical jump between 3 and 4
  clusters. What does this suggest?
\item
  Your average silhouette width is 0.25. Is this good clustering?
\item
  K-means gives you compact, equal-sized clusters, but hierarchical
  clustering with average linkage gives very different results. What
  should you do?
\end{enumerate}

\begin{center}\rule{0.5\linewidth}{0.5pt}\end{center}

\textbf{ }

\section{Answers to Practice Questions}

\textbf{ }

\subsection{Conceptual Answers}

\begin{enumerate}
\item
  \textbf{Cluster analysis} discovers unknown groupings (unsupervised),
  while \textbf{discriminant analysis} classifies into predefined groups
  (supervised).
\item
  Standardization ensures all variables contribute equally to distance.
  Example: Without standardization, income (0-200,000) would dominate
  over age (0-100) in Euclidean distance calculations.
\item
  Clusters form long chains rather than compact groups when similar
  observations link via intermediate points. \textbf{Single linkage} is
  most susceptible.
\item
  \textbf{Advantages:} Fast, scalable, simple, works well with large
  data. \textbf{Disadvantages:} Requires knowing k, sensitive to
  initialization, assumes spherical clusters, sensitive to outliers.
\item
  Silhouette close to +1 means the point is much closer to its own
  cluster than to neighboring clusters (well-clustered). Close to 0
  means borderline between clusters. Negative means likely in wrong
  cluster.
\end{enumerate}

\textbf{ }

\subsection{Application Answers}

\begin{enumerate}
\setcounter{enumi}{5}
\item
  \textbf{Yes, standardize.} Income (thousands) would dominate age,
  purchases, and loyalty\_score if not standardized, making the
  clustering essentially based only on income.
\item
  Use alternative methods: (1) Silhouette analysis, (2) Gap statistic,
  (3) Domain knowledge about appropriate number of segments, (4) Try
  multiple k values and compare interpretability.
\item
  Different \textbf{initializations} lead to different local optima.
  Solutions: (1) Run multiple times and choose best result (lowest
  WCSS), (2) Use k-means++ initialization, (3) Compare stability across
  runs.
\item
  \textbf{K-means} - hierarchical clustering would require O(\(n^{2}\))
  distance matrix (50,000 × 50,000 = 2.5 billion calculations). K-means
  scales much better to large n.
\item
  Not necessarily a problem. Small cluster might be: (1) Valuable niche
  segment (VIP customers), (2) Outliers that should be investigated, or
  (3) Indication that k is too high. Examine cluster characteristics to
  determine.
\end{enumerate}

\textbf{ }

\subsection{Interpretation Answers}

\begin{enumerate}
\setcounter{enumi}{10}
\item
  The point is \textbf{likely in the wrong cluster} - it's closer to a
  neighboring cluster than to its assigned cluster. Consider it
  misclassified.
\item
  \textbf{k=3 is better.} Lower Davies-Bouldin Index indicates more
  compact, better-separated clusters (0.5 \textless{} 0.8).
\item
  Suggests that \textbf{k=3 is optimal} - merging from 3 to 4 clusters
  causes a large increase in within-cluster heterogeneity, indicating
  natural separation at 3 clusters.
\item
  \textbf{Moderate clustering.} Silhouette width of 0.25 indicates some
  structure but not strong separation. Values 0.2-0.5 suggest reasonable
  but not excellent clustering. Consider if this is acceptable for your
  application.
\item
  \textbf{Compare both results carefully:} (1) Examine cluster profiles
  and interpretability, (2) Test stability with bootstrap samples, (3)
  Validate with domain knowledge, (4) Try intermediate methods like
  k-medoids. Different methods capture different aspects of structure -
  neither is necessarily ``wrong.''
\end{enumerate}

\begin{center}\rule{0.5\linewidth}{0.5pt}\end{center}

\textbf{ }

\section{Additional Resources}

\textbf{ }

\subsection{Recommended Reading}

\textbf{Textbooks:}

\begin{itemize}
\item
  Everitt, B., et al. (2011). \emph{Cluster Analysis} (5th ed.). Wiley.
\item
  Kaufman, L., \& Rousseeuw, P. J. (2005). \emph{Finding Groups in
  Data}. Wiley.
\item
  James, G., et al. (2021). \emph{An Introduction to Statistical
  Learning} (2nd ed.). Springer. Chapter 12.
\end{itemize}

\textbf{Papers:}

\begin{itemize}
\item
  Arthur, D., \& Vassilvitskii, S. (2007). k-means++: The advantages of
  careful seeding. \emph{SODA `07}.
\item
  Rousseeuw, P. J. (1987). Silhouettes: A graphical aid to the
  interpretation. \emph{Journal of Computational and Applied
  Mathematics}, 20, 53-65.
\end{itemize}

\textbf{ }

\subsection{Software Implementations}

\textbf{Python:}

\begin{itemize}
\item
  scikit-learn: KMeans, AgglomerativeClustering, DBSCAN
\item
  scipy.cluster: Hierarchical clustering and dendrograms
\item
  sklearn.metrics: Silhouette, Davies-Bouldin scores
\end{itemize}

\textbf{R:}

\begin{itemize}
\item
  stats: kmeans(), hclust(), dist()
\item
  cluster: pam(), silhouette()
\item
  factoextra: Visualization tools
\end{itemize}

\textbf{ }

\subsection{Online Resources}

\begin{itemize}
\item
  Scikit-learn Clustering Documentation: comprehensive guide with
  examples
\item
  StatQuest YouTube: Visual explanations of clustering concepts
\item
  Towards Data Science: Practical tutorials and case studies
\end{itemize}

\textbf{End of Lecture Notes}

Good luck with your studies!

\end{document}

